\chapter{化学键不是原子之间的小棍}

\begin{kaichang}
找一盒火柴，取出一根，划着它。

先别急着吹灭，认真看三秒钟：木棍头上的药料本来安安静静，空气里的氧气也安安静静，谁都没惹谁。可你只是在火柴盒侧面擦了一下，“砰”地一声轻响，火就来了——光来了，热也来了，烫得你得甩手。

现在猜一个问题：这份热，这份光，这份烫，\textbf{是谁付的账}？摩擦吗？可你擦完之后早就不擦了，火苗还在自己烧，还在持续放热。木头吗？可没点着的木头摸上去是凉的。空气吗？空气此刻正包围着你，你并没有着火。

把你的猜测写下来。这一章结束时，我们回到这根火柴，你会看到：那团火焰里的每一分能量，都是亿万个原子“结伴”时集体结的账。
\end{kaichang}

\section{看得见的现象：掰开费劲，结合发热}

先不急着钻到原子里面去。在肉眼的世界里，“原子结伴”这件事有两张面孔，你其实都见过。

第一张面孔：\textbf{把连在一起的东西分开，必须花力气}。一块饼干你能掰开，一根粉笔你也能掰断，但要掰断一根铁丝，徒手几乎不可能；想拉断一根棉线，得用两只手使劲；就算只是把两张粘在一起的便利贴撕开，你也得付出一点力。分开“结伴”的东西，从来都不是免费的。

第二张面孔：\textbf{东西结合起来的时候，常常放出能量}。火柴燃烧发光发热；氢气和氧气混合后点火，会“轰”地一声放出大量的热（所以氢气才被拿来当火箭燃料）；把两块磁铁靠近，它们会“啪”地吸在一起，撞上时发出声音——那声音也是能量。结合，往往是能量向外走的时候。

把这两张面孔放在一起，一条线索就浮现了：分开要付钱，结合会找零。自然界里仿佛有一本账簿，专门记录“原子结伴与分手”的收支。这一章的任务，就是把这本账簿翻开，看懂它的记账规则。

\section{把镜头推近：两个原子靠近时，发生了什么}

从第 9 章我们知道，原子是一个很小的、带正电的原子核，加上一团按能量分层居住的电子。现在，想象两个各自完整的原子，从很远的地方慢慢靠近。为了方便，就用最简单的氢原子：一个质子，一个电子。

两个原子离得很远时，互不搭理。可当它们近到电子“云层”开始交叠，\textbf{好几股力会同时变化}，而且方向并不一致：

\begin{itemize}[itemsep=2pt]
\item 甲原子的原子核开始吸引乙原子的电子，乙原子的核也吸引甲的电子。正负相吸，这是把两个原子往一起拉的力量。
\item 与此同时，两个带正电的原子核互相排斥，两团带负电的电子也互相排斥。同种电荷相斥，这是把两个原子往外推的力量。
\end{itemize}

注意，这不是“先吸引后排斥”的两幕戏，而是\textbf{同一时刻、所有粒子之间所有吸引与排斥的总和}。两个原子最终停在什么距离上、要不要结伴，不取决于任何单一的一股力，而取决于所有力加在一起之后，整个体系的能量是升了还是降了。

这就引出本章的核心判据，请把它划下来：\textbf{原子会不会成键，看的不是“谁喜欢谁”，而是整个体系（所有核、所有电子）的总能量。成键，就是体系找到了一个能量更低的摆法。}

原子没有愿望，不“想要”稳定，也不“为了凑八个电子”。它们只是遵循一条朴素得近乎乏味的规则：能量能从高往低走，就往低处走——和水往低处流是同一副脾气。

\section{一条曲线讲清成键：能量有一个谷底}

把“两个原子的距离”和“体系的总能量”画在一张图上，你会得到一条形状很像“山谷”的曲线。它大概是这一篇里最重要的一张图：

\begin{center}
\begin{tikzpicture}[scale=1.0]
\draw[->] (0,0) -- (8.2,0) node[right]{两个原子核之间的距离};
\draw[->] (0,-1.6) -- (0,3.0) node[above]{体系能量};
\draw[thick] plot[smooth] coordinates {(0.55,2.7) (0.8,1.1) (1.2,-0.3) (1.8,-0.9) (2.3,-1.0) (3.0,-0.72) (4.0,-0.34) (5.5,-0.12) (7.5,-0.03)};
\draw[dashed] (2.3,-1.0) -- (2.3,0);
\draw[<->] (2.75,-1.0) -- (2.75,0);
\node at (3.7,-0.5) {键能（谷有多深）};
\node at (2.3,0.35) {键长};
\node at (6.2,0.5) {相距很远：互不理睬};
\node at (1.0,2.2) {挤太近：};
\node at (1.0,1.75) {排斥猛增};
\end{tikzpicture}
\end{center}

（这是人为画的示意图，横轴纵轴都不对应某个具体分子的真实刻度，只表示趋势。）

从右往左读这条曲线，正好是“两个原子互相靠近”的过程：

\begin{itemize}[itemsep=2pt]
\item \textbf{离得很远}（图的右边）：谁也不影响谁，体系能量就取为零点，当作记账的基准。
\item \textbf{逐渐靠近}：核对电子的吸引开始占上风，总能量\textbf{下降}。就像水顺着山坡往下流，两个原子会自发地靠拢——如果有多余的能量能被带走的话。
\item \textbf{到达谷底}：在某个特定的距离上，吸引和排斥的账算得最划算，能量最低。这个距离就叫\textbf{键长}；从谷底回到零点所需要的能量（也就是谷的深度）就叫\textbf{键能}。
\item \textbf{再往里挤}：两个核之间的排斥、电子云之间的排斥开始暴涨，总能量掉头向上，曲线陡得几乎竖起来。原子不是可以随便压缩的软柿子。
\end{itemize}

所以“化学键”到底是什么？它\textbf{不是}一根把两个原子拴住的小棍，而是这样一个事实：在某个距离上，这个体系的能量恰好坐在谷底。想让原子待在谷底外面，你就得持续付钱（做功）；原子一旦掉进谷底，它自己就待在那儿了。

\begin{qiaoliang}
键能有多大？拿最真实的数字感受一下。氢气分子里那根 \chem{H-H} 键，键能约为 \ud{436}{kJ\cdot mol^{-1}}：拆散 \ud{1}{mol}（也就是 \ud{2}{g}）氢分子的键，必须支付 436 千焦的能量。平摊到单个键上：

\[ \frac{436\times 10^{3}\ \mathrm{J}}{\ud{6.02\times 10^{23}}{mol^{-1}}} \approx 7.2\times 10^{-19}\ \mathrm{J} \]

一根键的能量只有不到十亿亿分之一焦耳，小得可怜。可一枚再普通不过的火柴，烧完能放出约 \ud{1000}{J} 的热——这相当于大约 $10^{21}$ 个（一亿亿亿个）新键形成时一起“结账”放出的能量。单个原子的账微不足道，可一旦以摩尔为单位集体行动，就是能烫伤你手指的宏观火焰。这就是第 5 章认识过的阿伏伽德罗常数的威力：它把粒子世界的零钱，兑换成了我们世界的大钞。
\end{qiaoliang}

\section{能量从哪来、到哪去：断键吸能，成键放能}

现在把曲线翻译成记账语言，就是本章必须讲准的一句话：

\textbf{断开化学键，必须吸收能量；形成化学键，必然释放能量。}

方向千万不要记反。断键就是把原子从谷底往外拉，逆着吸引走，必须做功，就像把吸在一起的两块磁铁掰开要费劲一样。成键则是原子落进谷底，多余的能量以热、光等形式甩给外界，就像两块磁铁“啪”地吸在一起时发出声响。

\begin{shiyan}
\textbf{拉断与吸合：用磁铁和回形针体验“键的账簿”}

材料：两块能吸在一起的小磁铁（冰箱贴磁铁即可）、一个金属回形针、一根干意大利面或粉笔。

步骤：① 把两块磁铁吸在一起，再用力把它们拉开，体会手上必须付出的力。② 让它们从你手里“啪”地吸回去，凑近听那声轻响，并摸摸碰撞的地方。③ 用手把回形针来回弯折，折几次后立刻摸摸弯折处——是热的；继续弯折直到它断，体会断开之前你必须持续做功。④ 掰断干意面或粉笔，注意“咔嚓”那一下：断口出现之前，你的肌肉一直在付钱。

观察点：分开要结合的东西永远要花力气（吸能）；结合与碰撞会放出声音和热（放能）。回形针弯折处发热说明：你做的功没有消失，变成了热。

安全提示：只用小磁铁，防止夹手；弯折回形针时不要对着自己或他人的眼睛，断口锋利，做完洗手。这个实验里磁铁和金属的力是宏观的电磁力与金属内部的结合，只是帮助你体会“能量方向”的类比，细节第 18、21 章会讲。
\end{shiyan}

有了这条规则，再看化学反应就通透了。任何反应里，旧键要断开（吸能），新键要形成（放能），\textbf{反应整体放不放热，取决于两边轧账的结果}。以氢气在氧气中燃烧为例，用后面马上要见的键能表来算：

\[ \chem{2H_2} + \chem{O_2} \rightarrow \chem{2H_2O} \]

断键要支付：2 份 \chem{H-H}（$2\times 436$）加 1 份 \chem{O=O}（498），共约 \ud{1370}{kJ}；成键能收回：4 份 \chem{O-H}（$4\times 463$），共约 \ud{1852}{kJ}。收大于支，净放出约 \ud{482}{kJ}——这还只是按气态水估算的粗略值，真实燃烧放热更多。火箭之所以能把那么重的家伙推上天，靠的就是这笔“成键找零”。

反过来，把水分解回氢气和氧气，每一笔账都要倒着付——所以电解水必须持续通电，一分钱能量都省不掉。世界上不存在“把水轻轻一碰就变成燃料”的机器，谁宣称有，谁就是在否认这本账簿。

\section{符号层：那条小短线，只是记账符号}

到这儿，我们才请符号登场。你在课本上见过这样的写法：

\[ \chem{H-H} \qquad \chem{O=O} \qquad \chem{N\equiv N} \]

原子之间的小短线，一条叫单键、两条叫双键、三条叫三键（细节第 19 章会细讲）。现在你知道了：这些线\textbf{不是}画给原子用的连接件，而是人类记账用的符号——每一道线都代表“这里有一个能量谷底，大约这么深”。

下面是几种常见键的键能（数值是多种分子里测出来的平均值，同一个 \chem{C-H} 在甲烷里和在酒精里会略有出入）：

\begin{table}[h]
\centering
\begin{tabular}{lclc}
\toprule
键 & 键能（$\mathrm{kJ\cdot mol^{-1}}$） & 键 & 键能（$\mathrm{kJ\cdot mol^{-1}}$） \\
\midrule
\chem{H-H} & 436 & \chem{C-H} & 413 \\
\chem{C-C} & 347 & \chem{O-H} & 463 \\
\chem{O=O} & 498 & \chem{H-Cl} & 431 \\
\chem{N\equiv N} & 946 & \chem{Cl-Cl} & 243 \\
\bottomrule
\end{tabular}
\end{table}

扫一眼这张表，你能立刻看懂好几件日常小事。氮气里的 \chem{N\equiv N} 键能高达 \ud{946}{kJ\cdot mol^{-1}}，几乎是表里最深的谷底——所以空气里五分之四的氮气整天从你的肺里进进出出却几乎不参与任何反应，温顺得像个旁观者；而要把氮气“拆开”变成庄稼能用的氮肥，工厂不得不用高温高压硬砸这个谷底（这就是著名的哈伯法，第 30 章讲反应条件时会再遇到它）。氧气就活泼得多：\chem{O=O} 的键能只有氮气的一半左右，谷底浅，容易被“请”出来参加反应——这正是你能靠呼吸活着、火能烧得起来的化学原因之一。

键能越大，谷底越深，原子结伴越牢固；谷底越浅，这段关系越容易散伙。一张表，半本化学。

\section{八隅体：好用的口诀，不是法律}

第 10 章我们拆穿过一个流传极广的说法：“原子反应是为了凑够 8 个电子”。现在有了能量语言，可以把这件事彻底讲准。

“外层凑满 8 个电子就稳定”这条口诀（八隅体规则）确实好用：主族元素组成的绝大多数常见分子——水、二氧化碳、甲烷、食盐里的离子——都符合它。原因也很正当：对第二周期的碳、氮、氧、氟这类原子来说，外层住满 8 个电子的排布，往往\textbf{恰好}对应体系能量最低的谷底。口诀是对能量规律的一种速记。

但速记不是法律。只要抬出“总能量最低”这条真正的判据，例外立刻排着队出现：

\begin{itemize}[itemsep=2pt]
\item \textbf{不够 8 个也稳定}：三氟化硼 \chem{BF_3} 分子里，硼原子外层只有 6 个电子，可这个分子安安稳稳地存在，是工厂里常用的试剂。硼没有去“努力凑八”，因为照当时的处境，6 个电子的排法能量已经够低。
\item \textbf{超过 8 个也稳定}：六氟化硫 \chem{SF_6} 里，硫原子周围围着 12 个电子；五氯化磷 \chem{PCl_5} 里磷周围是 10 个。第三周期往下的原子“楼层”更大，住得下更多客人。
\item \textbf{奇数电子也拦不住}：一氧化氮 \chem{NO} 的电子总数是奇数，无论如何分，都不可能让每个原子都凑成成对的 8 个。可 \chem{NO} 不但稳定存在，还是你血管里重要的信号分子，此刻正在你体内帮助调节血压。
\item \textbf{连“贵族”都被拉下了水}：稀有气体外层恰好住满，曾被叫作“惰性气体”，教科书一度断言它们绝不成键。可 1962 年化学家造出了 \chem{XePtF_6}，随后 \chem{XeF_4} 等氙的化合物陆续问世——只要体系能量能降得足够低，“住满”的楼层也照样可以重新装修。
\end{itemize}

所以正确的说法是：\textbf{八隅体是有用的经验口诀，能量才是背后的法律}。口诀能帮你快速猜中大多数常见分子的结构，但每当你遇到例外，别惊讶，也别替原子“惋惜没凑够”——原子的账簿里从来没有“八个”这一项，只有“总能量”这一项。第 19 章画路易斯结构时，你还会和这条口诀的适用范围反复打交道。

\section{科学家怎样知道：从“小棍”到能量曲线}

\begin{zhengju}
化学键的概念比化学键的解释早出生了半个多世纪。1858 年前后，化学家库珀和凯库勒已经开始用短线把原子连起来写结构式——那条小短线好用极了，能解释为什么碳能搭出千奇百怪的骨架。但没有人能回答：那条线到底是什么东西？是原子伸出的钩子吗？没人见过。

1916 年，美国化学家路易斯提出一个大胆的猜法：成键是两个原子\textbf{共享一对电子}。这个猜法能解释一大批分子的组成，八隅体规则就是从这儿来的。可它立刻遭到质问：两个带负电的电子互相排斥，凭什么凑在一起反而能把原子粘住？这听起来违反直觉，路易斯自己也给不出计算。

真正的答案要等量子力学登场。1927 年，就在薛定谔方程发表的第二年，两位在德国工作的年轻物理学家海特勒和伦敦做了一个在当时极其困难的计算：把两个氢原子的核和电子全都当作量子体系，求解“整个体系的能量随核间距离怎么变”。算出来的曲线出现了一个明确的最低点——这就是本章开头那条山谷曲线第一次从理论里长出来。他们预言的键长和键能与实验测得的数值大体吻合，更重要的是，计算还说明了一件事：\textbf{单纯的经典静电吸引永远算不出这样的谷底}，成键本质上是量子效应（粗略地说，共享让两个电子的活动空间变大，动能得以下降，见本章挑战层）。后来鲍林等人把这套思想推广到各种分子，“化学键=体系能量降低”才从猜法变成了可以计算、可以检验的理论。

实验这边也没闲着。布拉格父子用 X 射线照射晶体，从衍射图样反推出原子之间的距离——键长第一次被“量”了出来；而分子振动吸收的红外光频率，则透露了键的“弹簧”有多硬，间接称出了键能。理论算的谷底深度和实验称的结果对得上，这套图景才真正站稳。注意这个故事的节奏：先有记账符号（1858），再有猜法（1916），然后才有解释（1927）和独立测量（衍射与光谱）。科学不是天才的顿悟接力，而是符号、猜法、计算、测量互相追认了将近一百年。
\end{zhengju}

\section{这个模型什么时候会失灵}

势能曲线和键能表是极其实用的工具，但你有权知道它们的边界，免得日后被它们反咬一口。

第一，曲线图把原子核当成了位置确定的小球，这是经典的画法。真实情况是电子以概率云的形式弥散在核周围（第 9 章），键长只是出现概率最大的距离，原子其实一刻不停地在谷底附近振动，绝不会安静坐着。

第二，键能表给的是\textbf{平均值}。同一个 \chem{O-H} 键，在水里和在酒精里的“成色”并不完全一样；查表算反应热，算出来的是八九不离十的估算，不是精确账。

第三，“一道短线一根键”的记账法，本身就有一大类记不了的账。苯分子里 6 个碳之间的键，既不是单键也不是双键，而是均匀分布的“一又二分之一”；金属铜里根本没有“哪两个原子结成一对”这回事，电子是整块金属共同持有的（第 21 章）；分子与分子之间那些微弱却重要的吸引——水为什么会流动成液体、壁虎为什么能爬墙——根本不在这张键能表上（第 22 章）。遇到这些情况，小棍模型不是“算错了”，而是压根不在它的记账范围内。

\begin{wuqu}
“化学反应释放能量，是因为断开了化学键。”——这是化学里传播最广的错误之一，方向完全记反了。

反例随手就有：如果断键放能，拉断铁丝应该\textbf{放出}热量，可你弯断回形针时摸到的是弯折处发热、断口并不喷火，而且你必须持续用力才断得开；如果断键放能，水分子应该自己散架变成氢气和氧气，顺便发电，而不是反过来必须通电才能拆开。

能量只在一个地方被释放：\textbf{原子落进能量谷底、新键形成的时候}。燃料燃烧放热，烧的不是“旧键断开”，而是“新键找零”——生成物（如 \chem{CO_2} 和 \chem{H_2O}）的键比反应物的键更深更稳，轧账之后净赚的那部分，才是你感受到的热。以后遇到“高能键断裂放出能量”之类的说法（包括某些对 ATP 的流行解释，第 46 章会专门清算它），请自动在心里把账重算一遍。
\end{wuqu}

\section{回到那根火柴}

现在兑现承诺。那根火柴的账，可以这样一笔一笔地算：

你在火柴盒侧面擦的那一下，摩擦做功，把一小份能量交给了火柴头上的药料。这份能量的用途非常具体——\textbf{支付第一笔断键费用}。药料里和空气中的旧键被撬开，原子们离开了各自的谷底。

接下来才是主角登场：这些被“解放”的原子重新结伴，一头扎进更深的新谷底——碳与氧结合成 \chem{CO_2}，氢与氧结合成 \chem{H_2O}，新键形成，\textbf{成键放能}。新谷底比旧谷底深得多，找零的能量以光和热的形式涌出，烫到你的手指，照亮火苗，顺便又把邻近的药料加热到“付得起断键费”的程度——于是火焰自己维持了下去，直到药料烧完。

所以火的本质是：\textbf{一场自掏腰包启动、靠成键找零维持的链式结账}。你划火柴付出的摩擦功只是“点火费”，真正的燃料是化学键之间的能量落差。这也顺便回答了开场时的迷惑：凉木头不放热，因为它的原子安分地待在谷底里；空气不烧你，因为 \chem{O=O} 的谷底虽然不深，也得先有人付那笔“断键首付”。

顺便说说开场提到的另一件事：铁丝为什么掰不断？因为金属里的原子不是两两结伴，而是所有原子共享一大片“公共电子”，整块金属就是一个巨大的能量谷底——那是另一种记账方式，第 21 章专门讲它。

\section{继续深挖：这本账簿通向哪里}

“断键吸能、成键放能”这条规则，远不只是解释火柴。它是整个能源世界的总账房：煤、石油、天然气之所以叫燃料，是因为它们的化学键经过重组能净放出能量；电池是把这笔账改成用电流的形式慢慢支付；炸药和普通燃料的账面完全一样，区别只在结账的\textbf{速度}——那是动力学的话题，第 28 章会见分晓。连你身体里也一样：你吃下的食物、呼吸的氧气，最终都要在这本账簿上轧账，为体温、心跳和思考付款——只不过生命结账的方式精巧得多，要等到第三册才慢慢展开。

不过，这一章我们只解决了“原子为什么结伴”的总原理，还没细看结伴的几种完全不同的方式。下一章先去围观最“干脆利落”的一种：钠遇到氯气时，电子不是被共享，而是\textbf{整个过户}——钠原子彻底失去一个电子，氯原子彻底抢走它，然后两种带电的离子靠静电引力抱成一团。这种“先转账、再拥抱”的成键方式，能解释为什么食盐坚硬又怕水、熔点高得能当坩埚材料，还能解释它为什么能点亮一盏灯。翻开第 18 章，看电子过户之后发生了什么。

\begin{tiaozhan}
（跳过不影响主线。）为什么共享电子能让能量降低？粗略的图像藏在海特勒—伦敦的计算里。量子力学中，电子的能量分两部分：被核吸引而降低的势能，以及与“活动空间大小”有关的动能——电子被限制得越死，动能越高，像被关在越小房间里的乒乓球跳得越凶。两个氢原子靠近时，每个电子都可以同时在两个核附近出现：活动空间变大了，动能下降；离两个正核都近了，势能也下降。两笔进账合起来压过了核与核的排斥，谷底就此形成。更精细地说，成键过程中动能与势能此消彼长，最终满足一条叫“维里定理”的严格关系。曲线谷底附近的一小段还可以近似成弹簧（抛物线），用 $E(r)=D\left(1-e^{-a(r-r_0)}\right)^2$ 这样的“莫尔斯势”来拟合整条曲线——分子的红外光谱正是通过测量这根“弹簧”的振动频率，反推出键有多硬。
\end{tiaozhan}

\section*{练一练}

\begin{sikaoti}
\item 不看课本，自己画一张“两个原子靠近时体系能量随距离变化”的曲线图，在图上标出：相距无穷远的零点、键长、键能，以及“挤得太近”时曲线为什么翘起来。用一句话向同学解释：为什么原子会“自发”停在谷底。
\item 查本章键能表，算一算反应 $\chem{H_2} + \chem{Cl_2} \rightarrow \chem{2HCl}$ 是放热还是吸热：断键共需支付多少能量，成键共能收回多少，轧账结果如何？（提示：\chem{H-Cl} 键能 \ud{431}{kJ\cdot mol^{-1}}。）
\item 有人说：“炸药威力大，是因为炸药分子里存了很多高能的键，爆炸时这些键断开，把能量放了出来。”请指出这句话里的两处方向性错误，并用“谷底”的语言重新讲一遍炸药为什么放能。
\item 氮气 \chem{N_2} 占空气约五分之四，却很少参与反应；氧气 \chem{O_2} 只占约五分之一，却能支持燃烧。用键能表解释这两种气体的“性格差”，并说说这个差异对你能安全呼吸意味着什么。
\item 硫原子在 \chem{SF_6} 分子中周围有 12 个电子，硼原子在 \chem{BF_3} 中只有 6 个。请用“八隅体是口诀、能量才是法律”的思路写一段话，向一个坚信“原子必须凑满 8 个电子”的朋友解释这两个分子为什么能稳定存在。
\end{sikaoti}
